% =====================================================================
% Exemple complet — POLYCOPIÉ — thème `legacy`
%
% Même contenu que themes/ocots/poly.tex, thème `legacy' seul change : c'est la
% base de comparaison entre thèmes, voir ../ (les autres dossiers de
% examples/themes/).
%
% Compilation : make themes/legacy/poly (depuis examples/)
% =====================================================================
\documentclass[11pt,twoside]{ocots-book}
\usepackage[lang=fr, theme=legacy, solutions=end, math=analysis,
            institution={n7}, author={Prénom Nom}]{ocots}

\title{Le template \texttt{ocots}}
\subtitle{Tout ce qu'il sait faire, sur un support papier --- thème legacy}
\author{Prénom \textsc{Nom}}
\date{\today}

\makeindex

\begin{document}

\dominitoc%

% ---------------------------------------------------------------- front
\frontmatter%
\maketitle
\tableofcontents

% ---------------------------------------------------------------- main
\mainmatter%

\part{Les environnements}%
\label{part:env}

\ocotscollectsolutions% début de la collecte des corrigés

% Contenu partagé — les environnements à boîte et de type théorème.
%
% Ce fichier ne contient aucun préambule : il est inclus par plusieurs
% documents, qui l'affichent avec des thèmes différents. C'est la base de
% comparaison v1 / v2.

\lstset{language=[LaTeX]TeX}

\chapter{Boîtes et théorèmes}%
\label{chap:boxes}
\minitoc%

\begin{chapterintro}
    Les résultats (théorème, définition, proposition, corollaire, conjecture,
    lemme, théorème cité) partagent un compteur : deux résultats d'une même
    section portent donc des numéros distincts. Les exemples, les remarques et
    les exercices ont chacun leur propre compteur. Chaque famille a une variante
    étoilée, non numérotée. Toutes prennent une liste de clés optionnelle :
    `title=', `label=' et `note='. Une étiquette est posée telle quelle, sans
    préfixe ajouté par le template.
\end{chapterintro}

%---------------------------------------------------------------------------
\section{Les familles à cadre}%
\label{sec:framed}

\begin{lstlisting}
\begin{theorem}[title={Theoreme avec titre}, label=thm:cauchy]
    Un enonce. Cette boite porte l'etiquette brute \texttt{thm:cauchy} et se cite
    par \verb|\ref{thm:cauchy}| : theoreme~\ref{thm:cauchy}.
\end{theorem}
\end{lstlisting}
\begin{theorem}[title={Théorème avec titre}, label=thm:cauchy]
    Un énoncé. Cette boîte porte l'étiquette brute \texttt{thm:cauchy} et se cite
    par \verb|\ref{thm:cauchy}| : théorème~\ref{thm:cauchy}.
\end{theorem}

\begin{lstlisting}
\begin{theorem}
    Un enonce sans titre ni etiquette : la liste de cles est facultative.
\end{theorem}
\end{lstlisting}
\begin{theorem}
    Un énoncé sans titre ni étiquette : la liste de clés est facultative.
\end{theorem}

\begin{lstlisting}
\begin{definition}[title={Definition}, label=def:ouvert, note={utile}]
    Un enonce. Reference~: definition~\ref{def:ouvert}.
\end{definition}
\end{lstlisting}
\begin{definition}[title={Définition}, label=def:ouvert, note={utile}]
    Un énoncé. Référence~: définition~\ref{def:ouvert}.
\end{definition}

\begin{lstlisting}
\begin{proposition}
    Un enonce.
\end{proposition}
\end{lstlisting}
\begin{proposition}
    Un énoncé.
\end{proposition}

\begin{lstlisting}
\begin{corollary}[title={Corollaire}]
    Un enonce.
\end{corollary}
\end{lstlisting}
\begin{corollary}[title={Corollaire}]
    Un énoncé.
\end{corollary}

\begin{lstlisting}
\begin{conjecture}[title={Conjecture}]
    Un enonce. Cette famille venait du fork HDR et manquait au template de ce
    depot.
\end{conjecture}
\end{lstlisting}
\begin{conjecture}[title={Conjecture}]
    Un énoncé. Cette famille venait du fork HDR et manquait au template de ce
    dépôt.
\end{conjecture}

%---------------------------------------------------------------------------
\section{Les environnements de type théorème}%
\label{sec:theoremlike}

Ceux-là ne sont pas encadrés : ils se composent au fil du texte, intitulé en
gras suivi du texte sur la même ligne.

\begin{lstlisting}
\begin{lemma}
    Un lemme numerote dans la section.
\end{lemma}
\end{lstlisting}
\begin{lemma}
    Un lemme numéroté dans la section.
\end{lemma}

\begin{lstlisting}
\begin{lemma*}
    Un lemme non numerote --- variante etoilee.
\end{lemma*}
\end{lstlisting}
\begin{lemma*}
    Un lemme non numéroté --- variante étoilée.
\end{lemma*}

\begin{lstlisting}
\begin{example}
    Un exemple. Il se termine par un carre blanc.
\end{example}
\end{lstlisting}
\begin{example}
    Un exemple. Il se termine par un carré blanc.
\end{example}

\begin{lstlisting}
\begin{example}[title={Cas $E=\R$}]
    Un exemple dont le titre contient une macro mathematique --- garde-fou de
    la regression de l'issue \#18 : la detection \texttt{title=} forcait
    autrefois l'expansion complete du titre (\texttt{xstring}) et cassait des
    qu'une macro non trivialement expansible y apparaissait, meme standard
    (\texttt{\textbackslash R}, \verb|\mathbb{R}|...).
\end{example}
\end{lstlisting}
\begin{example}[title={Cas $E=\R$}]
    Un exemple dont le titre contient une macro mathématique --- garde-fou de
    la régression de l'issue \#18 : la détection \texttt{title=} forçait
    autrefois l'expansion complète du titre (\texttt{xstring}) et cassait dès
    qu'une macro non trivialement expansible y apparaissait, même standard
    (\texttt{\textbackslash R}, \verb|\mathbb{R}|...).
\end{example}

\begin{lstlisting}
\begin{example*}[note={avec une note}]
    Un exemple non numerote, avec une note entre parentheses.
\end{example*}
\end{lstlisting}
\begin{example*}[note={avec une note}]
    Un exemple non numéroté, avec une note entre parenthèses.
\end{example*}

\begin{lstlisting}
\begin{remark}
    Une remarque, marquee par un filet lateral.
\end{remark}
\end{lstlisting}
\begin{remark}
    Une remarque, marquée par un filet latéral.
\end{remark}

\begin{lstlisting}
\begin{remark*}
    Une remarque non numerotee.
\end{remark*}
\end{lstlisting}
\begin{remark*}
    Une remarque non numérotée.
\end{remark*}

\begin{lstlisting}
\begin{remark*}[Titre brut]
    Une variante etoilee avec un titre brut, sans `='. C'est la syntaxe
    encore en usage dans le corpus (par exemple \texttt{td/td3.tex}) : elle
    doit continuer a fonctionner, pas seulement la forme \texttt{title=}.
\end{remark*}
\end{lstlisting}
\begin{remark*}[Titre brut]
    Une variante étoilée avec un titre brut, sans `='. C'est la syntaxe
    encore en usage dans le corpus (par exemple \texttt{td/td3.tex}) : elle
    doit continuer à fonctionner, pas seulement la forme \texttt{title=}.
\end{remark*}

%---------------------------------------------------------------------------
\section{Compteurs séparés et variantes étoilées}%
\label{sec:counters}

Les résultats partagent un compteur, les exemples et les remarques ont chacun
le leur. Dans la suite ci-dessous, la définition et la proposition sont donc
numérotées 1 et 2, quel que soit le nombre d'exemples ou de remarques placés
entre elles : des diapositives qui reprennent ce polycopié sans la remarque
gardent exactement les mêmes numéros de résultats.

\begin{lstlisting}
\begin{definition}[title={Ensemble negligeable}]
    Un enonce.
\end{definition}
\begin{remark}
    Une remarque : son numero vient de son propre compteur.
\end{remark}
\begin{proposition}
    Un enonce, numerote juste apres la definition.
\end{proposition}
\end{lstlisting}
\begin{definition}[title={Ensemble négligeable}]
    Un énoncé.
\end{definition}
\begin{remark}
    Une remarque : son numéro vient de son propre compteur.
\end{remark}
\begin{proposition}
    Un énoncé, numéroté juste après la définition.
\end{proposition}

Chaque famille a une variante étoilée : même habillage, même titre, mais ni
numéro ni incrément de compteur. Elle sert à reprendre, dans des diapositives,
une remarque ou un exemple sans perturber les numéros du polycopié, ou à
rappeler un résultat d'un autre cours. Une étiquette posée sur une variante
étoilée n'aurait rien à désigner : elle est ignorée, avec un avertissement.

\begin{lstlisting}
\begin{theorem*}[title={Pythagore}]
    Un enonce rappele, non numerote.
\end{theorem*}
\end{lstlisting}
\begin{theorem*}[title={Pythagore}]
    Un énoncé rappelé, non numéroté.
\end{theorem*}

\begin{lstlisting}
\begin{definition*}
    Une definition non numerotee.
\end{definition*}
\end{lstlisting}
\begin{definition*}
    Une définition non numérotée.
\end{definition*}

\begin{proposition*}[note={rappel}]
    Une proposition non numérotée, avec une note.
\end{proposition*}

\begin{corollary*}
    Un corollaire non numéroté.
\end{corollary*}

\begin{conjecture*}
    Une conjecture non numérotée.
\end{conjecture*}

\begin{citedtheorem*}
    Un théorème cité, non numéroté.
\end{citedtheorem*}

Les numéros reprennent ensuite là où ils en étaient : la proposition suivante
suit directement celle du début de la section.

\begin{proposition}
    Un énoncé numéroté, qui suit la proposition précédente.
\end{proposition}

%---------------------------------------------------------------------------
\section{Un même énoncé, deux supports}%
\label{sec:portable}

C'est le point de la refonte. L'énoncé ci-dessous est copié tel quel depuis
l'exemple de TD~: en v0 il aurait fallu le réécrire, parce que
\texttt{theorem} et \texttt{question} n'avaient pas la même signature d'un
support à l'autre.

\begin{lstlisting}
\begin{exercise}[nosolution]
    Soit $E$ un espace vectoriel norme.
    \begin{question}
        Montrer que la boule $\ball(0,1)$ est convexe.
    \end{question}
    \begin{question}
        En deduire que $\norm{\cdot}$ est convexe.
    \end{question}
\end{exercise}
\end{lstlisting}
\begin{exercise}[nosolution]
    Soit $E$ un espace vectoriel normé.
    \begin{question}
        Montrer que la boule $\ball(0,1)$ est convexe.
    \end{question}
    \begin{question}
        En déduire que $\norm{\cdot}$ est convexe.
    \end{question}
\end{exercise}

% Contenu partagé — preuves, blocs étiquetés, mises en valeur.

\chapter{Preuves et blocs annexes}%
\label{chap:blocks}
\minitoc%

%---------------------------------------------------------------------------
\section{Preuves}%
\label{sec:proofs}

\begin{proposition}{Une proposition à démontrer}{ademontrer}
    Pour tout $x \in \R$, $x^2 \ge 0$.
\end{proposition}

\begin{proof}
    Le marqueur d'ouverture et le carré de fin viennent du thème, comme la
    couleur du texte. \newstep{} Un astérisque sépare deux étapes.
\end{proof}

\begin{proof}[Démonstration de la proposition~\ref{prop:ademontrer}]
    L'argument facultatif remplace le marqueur d'ouverture.
\end{proof}

Une preuve peut être étalée sur plusieurs diapositives. Sur papier, les trois
morceaux se suivent~; en diapositives, seul le dernier pose le carré final.

\begin{proofbegin}
    Premier morceau.
\end{proofbegin}
\begin{proofmiddle}
    Morceau intermédiaire.
\end{proofmiddle}
\begin{proofend}
    Dernier morceau.
\end{proofend}

%---------------------------------------------------------------------------
\section{Blocs étiquetés}%
\label{sec:tagged}

Trois familles numérotées indépendamment, marquées par un filet latéral ou un
simple retrait. Elles partagent la même fabrique interne~: en ajouter une
quatrième tient en une ligne.

\begin{assumption}
    Une hypothèse. Elles sont étiquetées H1, H2, \ldots
\end{assumption}

\begin{assumption}
    Une deuxième hypothèse.
\end{assumption}

\begin{assumption*}
    Une hypothèse non étiquetée.
\end{assumption*}

\begin{openquestion}
    Un questionnement à se poser. Étiquetage Q1, Q2, \ldots
\end{openquestion}

\begin{difficulty}
    Une difficulté attendue. Étiquetage D1, D2, \ldots
\end{difficulty}

\begin{web}[Documentation en ligne du cours]
    Renvoi vers une ressource extérieure, précédé de son pictogramme.
\end{web}

%---------------------------------------------------------------------------
\section{Mise en valeur}%
\label{sec:emphasis}

Le \keyword{terme important} se compose avec \verb|\keyword|. Les quatre
couleurs du thème sont accessibles~: \emphA{première}, \emphB{deuxième},
\emphC{troisième}, \emphD{quatrième}, plus \emphLink{la couleur des liens}.

Pictogrammes~: \cmark{} et \xmark{}, plus la main \handwrite{} qui précédait les
exercices en v0.

\supplement{Un complément, en vert.}

\reminder{Un rappel, en rouge.}

Une équation mise en valeur par un surlignage, dont la couleur vient du thème~:
\[
    \tcbhighmath{\dot{x}(t) = A\, x(t) + b(t)}
\]

Numérotation indépendante des problèmes~:
\begin{equation}
    \min_{u} \int_0^{t_f} \! \ell(x(t), u(t)) \dif t
    \tagProblem
\end{equation}

%---------------------------------------------------------------------------
\section{Notes de travail}%
\label{sec:draft}

Visibles seulement avec l'option \texttt{draft}, invisibles sinon. Elles
remplacent le \verb|\ifVIEWALL| de la v0, qui n'était jamais défini et
provoquait une erreur à l'usage.

\notework{Une note personnelle.}

\worktodo{Quelque chose à finir.}

\begin{worknotes}
    Un bloc entier de notes de travail, exclu de la compilation finale.
\end{worknotes}

% Contenu partagé — exercices et corrigés.
%
% Le sort des corrigés dépend uniquement de l'option solutions= du paquet.
% Ce fichier est identique dans les trois modes.

\chapter{Exercices et corrigés}%
\label{chap:exercises}
\minitoc%

\begin{quotation}
    L'option \texttt{solutions=none} ne compose aucun corrigé --- c'est la
    version étudiante. \texttt{solutions=inline} les place sous l'énoncé.
    \texttt{solutions=end} les reporte en fin de partie, avec un renvoi dans
    chaque sens. Le fichier source, lui, ne change pas.
\end{quotation}

%---------------------------------------------------------------------------
\section{Formes d'exercice}%
\label{sec:exforms}

\begin{exercise}
    Un exercice simple, sans étiquette ni barème.
\solution
    Son corrigé.
\end{exercise}

\begin{exercise}[label=ex:matrices, points=4]
    Un exercice étiqueté et barémé. On le cite par \verb|\ref{ex:matrices}|.
    Calculer $e^{tA}$ pour $A = \begin{pmatrix} 0 & 1 \\ -1 & 0\end{pmatrix}$.
\solution
    On part de la définition $e^{tA} = \sum_{n \ge 0} \frac{t^n}{n!} A^n$ et de
    l'identité $A^2 = -I$, d'où
    \[
        e^{tA} = \begin{pmatrix} \cos t & \sin t \\ -\sin t & \cos t\end{pmatrix}.
    \]
\end{exercise}

L'exercice précédent est l'exercice~\ref{ex:matrices}.

\begin{exercise}[nosolution]
    Un exercice sans corrigé, quel que soit le mode~: la clef
    \texttt{nosolution} l'exclut individuellement, et le titre ne porte alors
    aucun renvoi.
\end{exercise}

%---------------------------------------------------------------------------
\section{Questions et sous-questions}%
\label{sec:questions}

\begin{exercise}[points=6]
    \begin{question}
        Une première question, numérotée dans l'exercice.
    \end{question}
    \begin{question}
        Une deuxième question.
        \begin{subquestion}
            Une sous-question.
        \end{subquestion}
        \begin{subquestion}
            Une autre sous-question.
        \end{subquestion}
    \end{question}
    \begin{question}
        Une troisième question.
    \end{question}
\solution
    Le compteur de questions repart à zéro dans le corrigé~; \verb|\newquestion|
    permet de le faire avancer à la main quand le corrigé ne suit pas l'énoncé
    question par question.
    \begin{question}
        Corrigé de la première question.
    \end{question}
    \begin{question}
        Corrigé de la deuxième.
    \end{question}
\end{exercise}

%---------------------------------------------------------------------------
\section{Correction hors boîte}%
\label{sec:correction}

Dans un TD ou un sujet d'examen, l'énoncé n'est pas toujours encadré. Le bloc
\texttt{correction} suit alors le texte et disparaît en mode
\texttt{solutions=none}.

\begin{correction}
    Une correction brève, au fil du texte.
\end{correction}


\chapter{Corrections des exercices}%
\label{chap:solutions}
\ocotsprintsolutions% les corrigés collectés viennent ici

% ------------------------------------------------------------------
\part{Les macros}%
\label{part:macros}

% Contenu partagé — les macros mathématiques.
%
% Ce chapitre exerce math=base et math=analysis. Les modules `control' et
% `measure' ne sont pas chargés ici : c'est justement l'intérêt de l'option,
% un polycopié de calcul différentiel n'a pas à embarquer les notations du
% contrôle optimal.

\chapter{Macros mathématiques}%
\label{chap:math}
\minitoc%

%---------------------------------------------------------------------------
\section{Ensembles de nombres}%
\label{sec:numbers}

Les ensembles usuels~: $\N$, $\Z$, $\Q$, $\R$, $\C$, $\K$, $\Sn$.

Leurs décorations, construites par \verb|\setDeco|~:
$\Rp$, $\Rn$, $\Rs$, $\Rsp$, $\Rsn$, $\Rb$, $\Rbp$, $\Ns$, $\Nb$.

Ensembles et espaces~: $\M_n(\R)$, $\B$, $\E$, $\F$, $\D$, $\O$, $\P$,
$\Vcal$, $\Ical$, $\Sgot$, et l'indicatrice $\ind_A$.

%---------------------------------------------------------------------------
\section{Normes, intervalles, opérateurs}%
\label{sec:norms}

Valeur absolue et norme, en version fixe ou extensible~:
$\abs{x}$, $\absStyle{\frac{a}{b}}$, $\norm{x}$, $\normStyle{\frac{a}{b}}$.

Produit scalaire $\prodscal{x}{y}$ et ensemble défini en compréhension~:
\[
    \enstq{x \in \R^n}{\norm{x} \le 1}.
\]

Intervalles~: $\intervalleff{a}{b}$, $\intervallefo{a}{b}$,
$\intervalleof{a}{b}$, $\intervalleoo{a}{b}$, et l'intervalle d'entiers
$\intervalleentier{1}{n}$.

\begin{remark}
    La convention des intervalles ouverts vient du fichier de langue~: $]a,b[$
    en français, $(a,b)$ en anglais. C'est ce qui évite de dupliquer tout le
    fichier de macros pour changer de langue.
\end{remark}

Opérateurs~: $\rang A$, $\rank A$, $\trace A$, $\det A$, $\Ker A$, $\im A$,
$\vect(u,v)$, $\diag(a_1,\ldots,a_n)$, $\card E$, $\supp f$, $\sign x$,
$\graphe f$, $\GL_n(\R)$, $\id$.

Les noms de $\rang$, $\minimize$ et $\graphe$ dépendent de la langue.

Comparaison de matrices~: $A \semidefpos$, $A \defpos$, $A \defneg$.

%---------------------------------------------------------------------------
\section{Applications et différentielles}%
\label{sec:diff}

Une application, avec \verb|\fonction|~:
\[
    \fonction{f}{\R^n}{\R^p}{x}{f(x)}
\]

Éléments différentiels~: $\xdif t$, $\xDif f(x)$, $\pardiff_x f$.

Dérivées partielles~: $\frp{f}{x}$, $\frpp{f}{x}$, $\frpij{f}{x}{y}$,
$\frpxx{f}$, $\frptt{f}$.

Gradient, avec ou sans indice~: $\grad{f}$ et $\grad{f}{x}$.

Notations de Landau~: $g(v) = \petito{v^2}$ et $h(v) = \grandO{v}$.

Espaces fonctionnels~: $\xCn{1}(\Ical, \R^n)$, $\xLn{2}(E,F)$, $\Ccal$,
$\Lcal$, $\Dcal$, $\Ucal$, $\Acal$, $\Ecal$, $\Fcal$, $\Kcal$, $\Ncal$,
$\GLcal$, $\Htrue$.

Topologie~: la boule ouverte $\Ball(x,r)$, la boule fermée $\BallClosed(x,r)$,
l'adhérence $\adherence{A}$, et la barre ajustée $\xoverline{AB}$ qui remplace
$\bar{\phantom{x}}$ sur les symboles larges.

Quotient~: $\EnsembleQuotient{E}{R}$. Convergence~: $u_n \convn u$.

%---------------------------------------------------------------------------
\section{Un énoncé complet}%
\label{sec:realstatement}

\begin{theorem}{Cauchy--Lipschitz}{cauchylipschitz}
    Soit $f \colon \Ical \times \Omega \to \R^n$ continue, localement
    lipschitzienne en $x$, avec $\Ical$ un intervalle ouvert de $\R$ et
    $\Omega$ un ouvert de $\R^n$. Alors pour tout
    $(t_0, x_0) \in \Ical \times \Omega$, le problème de Cauchy
    \[
        \dot{x}(t) = f(t, x(t)), \qquad x(t_0) = x_0,
    \]
    admet une unique solution maximale, définie sur un intervalle ouvert
    $\intervalleoo{t_-}{t_+} \subset \Ical$ contenant $t_0$.
\end{theorem}

\begin{proof}
    Le point fixe de Picard sur $\xCn{0}(\intervalleff{t_0-\eta}{t_0+\eta},
    \BallClosed(x_0,r))$, muni de $\norm{\cdot}_\infty$.
\end{proof}

% Contenu partagé — dessins, tableaux, code source.

\chapter{Dessins, tableaux et code}%
\label{chap:graphics}
\minitoc%

%---------------------------------------------------------------------------
\section{Dessins}%
\label{sec:tikz}

Les aides TikZ du template posent une fenêtre, une grille et des axes sans
avoir à les redessiner à chaque figure.

\begin{center}
\begin{tikzpicture}[xmin=-2, xmax=2, ymin=-1.5, ymax=1.5, scale=1.2]
    \grille
    \axes{$t$}{$x$}
    \draw[thick, color=\ocotscolor{emph-a}, domain=-2:2, samples=100]
        plot (\x, {sin(\x r)});
\end{tikzpicture}
\end{center}

Les pointes de flèches par défaut viennent du template~:
\begin{center}
\begin{tikzpicture}
    \draw[->] (0,0) -- (3,0) node[right]{$x$};
    \draw[->] (0,-0.3) -- (0,1) node[above]{$y$};
\end{tikzpicture}
\end{center}

%---------------------------------------------------------------------------
\section{Annotation d'une image}%
\label{sec:tikzgraphics}

L'environnement \texttt{tikzgraphics} permet de placer des repères sur une
image bitmap en coordonnées pixel, sans recalcul d'échelle.

\begin{center}
\begin{tikzgraphics}{6cm}{2546}{552}{logo-insa}
    \pxnode[anchor=west, color=\ocotscolor{emph-b}]{2600}{276}{repere}{$\leftarrow$ repère}
\end{tikzgraphics}
\end{center}

%---------------------------------------------------------------------------
\section{Tableaux}%
\label{sec:tables}

Colonnes centrées et alignées à gauche de largeur fixe, filets d'épaisseurs
graduées, cales verticales.

\begin{center}
\begin{tabular}{C{3cm} C{3cm} C{3cm}}
    \bighrule
    Méthode & Ordre & Stabilité \Tstrut\Bstrut \\
    \medhrule
    Euler explicite  & 1 & conditionnelle \Tstrut \\
    Euler implicite  & 1 & inconditionnelle \\
    Runge--Kutta 4   & 4 & conditionnelle \Bstrut \\
    \bighrule
\end{tabular}
\end{center}

%---------------------------------------------------------------------------
\section{Code source}%
\label{sec:listings}

Le style de \texttt{listings} suit les couleurs du thème.

\begin{lstlisting}[language=Python, caption={Une méthode d'Euler explicite}]
def euler(f, t0, x0, tf, n):
    # pas de discretisation
    h = (tf - t0) / n
    t, x = t0, x0
    for k in range(n):
        x = x + h * f(t, x)
        t = t + h
    return t, x
\end{lstlisting}

%---------------------------------------------------------------------------
\section{Liens, URL, renvois}%
\label{sec:links}

Un lien nommé vers une
\href{https://github.com/ocourses/ocots-latex-template}{page du dépôt}, une URL
nue \url{https://ocourses.github.io/}, une note de bas de page qui renvoie
ailleurs\footnote{Voir aussi \url{https://www.ctan.org/pkg/tcolorbox}.}, et
\emphLink{un fragment coloré à la main} avec \verb|\emphLink|.

Les renvois internes traversent les chapitres~: le théorème~\ref{thm:cauchy} du
chapitre~\ref{chap:boxes}, la définition~\ref{def:ouvert}, la
proposition~\ref{prop:ademontrer} du chapitre~\ref{chap:blocks}, et
l'exercice~\ref{ex:matrices} du chapitre~\ref{chap:exercises}.

Index~: le mot \index{flot}flot est indexé, ainsi que le
mot \index{orbite}orbite.


% ------------------------------------------------------------------
\part{Annexes}%
\label{part:appendix}

\begin{appendix}
    \chapter{Une annexe}%
    \label{chap:appendix}
    \minitoc%

    Les annexes sont numérotées en lettres.
\end{appendix}

% ---------------------------------------------------------------- back
\backmatter%
\printindex%

\end{document}
